\subsection{Infinite plays and objective-sensitive fixpoints}\label{app:infinite-plays}

Infinite plays require an objective for the infinite trace. They do not
require changing the public owned-choice operation when continuations
have equilibria. We separate this composition fact from a fixpoint
description of full subgame perfection.

The general equilibrium-existence and definability questions have
established answers. Ummels \cite[Theorems~6 and~13]{Ummels2006} proves
pure SPE existence for qualitative Borel multiplayer objectives and
modal $\mu$-calculus definability of constrained SPE existence when the
finite player set, the $\omega$-regular objectives and the payoff
constraints are fixed. The
construction below instead gives an explicit, self-contained formula
for a specified complete profile. Profiles remain quantified; we do not
use an existence or determinacy theorem in its proof.

\paragraph{Public composition.}
Write $E(G)$ for the possible winners of pure SPEs, allowing $E(G)$ to
be empty. Strategies retain decisions in every subgame, including
unreached ones. Each play has exactly one winner, and player $i$'s
payoff is the indicator of that winner being $i$. Each continuation
keeps its original winner objective when grafted.

\begin{proposition}[Public grafting]\label{olv:prop:infinite-grafting}
Let $G,H$ be continuation games, finite or infinite, with the stated
winner payoffs. Adding one public choice owned by $r$ gives
\[
 E(G\ow r H)=
 \begin{cases}
  \varnothing,&E(G)=\varnothing\text{ or }E(H)=\varnothing,\\
  \{r\},&\text{both values are nonempty and one is }\{r\},\\
  E(G)\cup E(H),&\text{otherwise}.
 \end{cases}
\]
In particular, on nonempty continuation values this is exactly the
paper's owned choice.
\end{proposition}
\begin{proof}
A full profile at the graft consists of the root choice and complete
profiles in both children. Subgame perfection requires both child
profiles to be SPEs, proving the empty case. Fix child SPEs with
winners $a,b$. The root profile is Nash precisely when its selected
winner is $r$ whenever $a=r$ or $b=r$. For necessity, the root owner
can switch branches without changing either continuation. For
sufficiency, a nonowner's deviation stays in the selected child and
is unprofitable by its Nash property. The owner's deviation can
select either child, but its Nash property prevents an improvement
there unless that child's original winner is already $r$.

If a child value is $\{r\}$ this condition forces $r$. Otherwise
each child has a losing SPE for $r$. Any desired non-$r$ winner in
either value is realized by selecting its SPE and taking such a
losing SPE in the other child. A desired winner $r$ is realized by
selecting an SPE that awards $r$. These are all possible winners.
\end{proof}

\paragraph{Finite-priority objectives.}
Let the public history tree be $M^*$: every node offers the same
nonempty move set $M$, while its owner and priority may depend on the entire history.
Assign priorities in a finite nonempty chain and assign each priority
one winning role. An infinite play is won by the role assigned to its
largest priority occurring infinitely often. Earlier finite prefixes
do not affect the winner of a subgame. This includes total
action-labelled cyclic arenas by unfolding. Finite terminating games
can be included by replacing a terminal winner by a constant-priority
absorbing continuation.

Fix a complete profile $\sigma$, a move at every history. Let
$R_\sigma$ follow that move. For player $i$, let $R_i^\sigma$ allow
every move at $i$'s nodes and follow $\sigma$ at other nodes. Write
$\operatorname{Pre}_R X=\{h:\exists k,\ hRk\land k\in X\}$.
All fixpoints below are on the complete lattice of sets of histories.
For a priority $p$, put $P_p=\{h:\operatorname{priority}(h)=p\}$ and
$Q_p=\{h:\operatorname{priority}(h)\leq p\}$, and define
\begin{align*}
 B_p^R&=\nu Z.\,\mu Y.\,
       (P_p\cap\operatorname{Pre}_R Z)
       \cup(Q_p\cap\operatorname{Pre}_R Y),\\
 W_i^R&=\bigcup_{p:\operatorname{winner}(p)=i}
       \mu X.\,(B_p^R\cup\operatorname{Pre}_R X),\\
 N_\sigma&=\{h:\forall i,\ h\in W_i^{R_i^\sigma}
                              \Longrightarrow h\in W_i^{R_\sigma}\},\\
 S_\sigma&=\nu Z.\,(N_\sigma\cap\square Z),
\end{align*}
where $\square Z$ requires every one-move continuation to be in $Z$.

\begin{theorem}[Full SPE by fixpoints]\label{olv:thm:infinite-fixpoints}
For every complete profile $\sigma$ and history $h$,
\[
 \sigma\text{ is subgame perfect from }h
 \quad\Longleftrightarrow\quad h\in S_\sigma.
\]
Consequently the possible SPE winners are exactly
\[
 E(G,h)=\{i:\exists\sigma,\ h\in S_\sigma\cap W_i^{R_\sigma}\}.
\]
The equivalence uses arbitrary history-dependent strategies and whole
unilateral replacements, with no equilibrium-existence premise.
\end{theorem}
\begin{proof}
First, for any relation $R$, $B_p^R$ is exactly the set of starts of
infinite $R$-paths staying in $Q_p$ and visiting $P_p$ arbitrarily late.
For the forward direction, unfold the outer greatest fixpoint. The
inner least fixpoint supplies a finite path to a $P_p$ state with
an edge back into the greatest fixpoint. At each state choose a
witness of minimum finite length. Before reaching length zero,
one step strictly decreases this length; at length zero take the
return edge. Iterating gives an infinite path, and a length-zero
state occurs after every finite prefix. Conversely, the states of
any such path form a postfixed set: from every position a later
$P_p$ visit followed by its next edge witnesses the inner least
fixpoint. Greatest-fixpoint coinduction gives membership.

The outer least fixpoint in $W_i^R$ permits a finite prefix before
this repeatedly visiting tail. An infinite sequence over a finite priority set has a
largest recurrent priority, and every larger priority occurs only
finitely often. Thus $h\in W_i^R$ means exactly that some infinite
$R$-path from $h$ is won by $i$.

For $R_\sigma$ the path is the profile's unique play. An
$R_i^\sigma$-path is exactly the play of some complete unilateral
replacement by $i$. In the reverse direction, prescribe its moves
on the path and use $\sigma$ everywhere else; path histories have
distinct lengths, so these prescriptions are consistent. Hence
$N_\sigma$ is exactly the set of histories where $\sigma$ is Nash
against complete deviations. Finally,
$\nu Z.(N_\sigma\cap\square Z)$ consists exactly of histories
whose every finite continuation belongs to $N_\sigma$. This is
the independent definition of subgame perfection. Existentially
quantifying profiles gives the winner formula.
\end{proof}

Whole deviations matter. Suppose $a$ repeatedly chooses Continue or
Exit; Exit awards $b$, while continuing forever awards $a$. The
always-Exit profile cannot be improved by changing its move at finitely
many histories, even after any history. The complete always-Continue
replacement improves it from the initial history. The path test above
detects this deviation; a one-step or finite-deviation test does not.

The least and greatest fixpoints concern the objective-sensitive
history predicates, not just sets of possible winners. Taking the
least fixpoint of $X\mapsto X\ow{a}^{0}\{a\}$, where
$\ow{a}^{0}$ is the strict empty-value extension, gives the empty set.
Yet in the game where $a$ chooses between repeating this choice and
exiting to its own win, with infinite repetition awarded to $b\ne a$,
the SPE value is $\{a\}$. Public grafting is therefore compatible
with infinite continuations, while recursive feedback needs the
trace objective as well as the finite winner observation.
